\section{Preliminaries}  \label{sec:prel}
Before proving the results, we describe some of the assumptions and notations used throughout the paper. $\x$ denotes $(x_1,\hdots,x_n)$. $[n]$ denotes $\{1,\hdots,n\}$.

\subsection{Notations and Definitions}

\smallskip
\begin{itemize}
    \item \textbf{Logarithmic derivative.} Over a ring $\ringR$ and a variable $y$, the logarithmic derivative $\dlg_{y} : \ringR(y) \to \ringR(y)$ is defined as $\dlg_{y}(f):=\partial_y\,f/f$; here $\partial_{y}$ denotes the partial derivative with respect to variable $y$. One important property of $\dlg$ is that it is additive over a product as \[\dlg_{y}(f \cdot g) \,=\, \frac{\partial_y(f \cdot g)}{f \cdot g} \,=\, \frac{(f \cdot \partial_y g \,+\, g \cdot \partial_y f)}{f \cdot g}\,=\,\dlg_{y}(f) + \dlg_y(g).\] 
    \noindent We refer this effect as \emph{linearization} of product.

    \item \textbf{Circuit size.} Sparsity $\sp(\cdot)$ refers to the number of nonzero monomials. In this paper, it is a parameter of the circuit size. In particular, $\size(g_1 \cdots g_s)=\sum_{i \in [s]}\,(\sp(g_i)+\deg(g_i))$, for $g_i \in \Sigma \wedge$ (respectively~$\Sigma \Pi^{[\delta]}$). In whitebox settings, we also include the \emph{bit-complexity} of the circuit (i.e.~bit complexity of the constants used in the wires) in the size parameter. 
    Some of the complexity parameters of a circuit are \emph{depth} (number of layers), \emph{syntactic degree} (the maximum degree polynomial computed
    by any node), {\em fanin} (maximum number of inputs to a node).

    \item \textbf{Hitting set.} A set of points $\calH \subseteq \F^n$ is called a \emph{hitting-set} for a class $\calC$ of $n$-variate polynomials if for any nonzero polynomial $f \in \calC$, there exists a point in $\calH$ where $f$ evaluates to a nonzero value. A $T(n)$-time hitting-set would mean that the hitting-set can be generated in time $T(n)$, for input size $n$. 

    \item \textbf{Valuation.} Valuation is a map $\val_y : \ringR[y] \to \Z_{\ge 0}$, over a ring $\ringR$, such that $\val_y(\cdot)$ is defined to be the maximum power of $y$ dividing the element. It can be easily extended to fraction field $\ringR(y)$, by defining $\val_{y}(p/q):= \val_y(p) - \val_{y}(q)$; where it can be negative.
    
    \item \textbf{Field.} We denote the underlying field as $\F$ and assume that it is of characteristic $0$. All our results hold for other fields (eg.~$\Q_p, \F_p$) of \emph{large} characteristic (see Remarks in Section \ref{sec:pf:thm1}-\ref{sec:pf-thm2}). 
    
    \item \textbf{Jacobian.} The Jacobian of a set of polynomials ${\mathbf f}=\{f_1,\hdots,f_m\}$ in $\F[\x]$ is defined to be the matrix $\calJ_{\x}({\mathbf f}) := \left(\partial_{x_j}(f_i)\right)_{m \times n}$. Let $S \subseteq \x = \{x_1,\hdots,x_n\}$ and $|S|=m$. Then, polynomial $J_{S}({\bf f})$ denotes the minor (i.e.~determinant
    of the submatrix) of $\calJ_{\x}({\bf f})$, formed by the columns corresponding to the variables in $S$.
\end{itemize}

\subsection{Basics of Algebraic Complexity Theory} \label{sec:basic-act}

For detailed discussion on the basics of Algebraic Complexity Theory we will encourage readers to refer \cite{shpilka2010arithmetic,saxena2009progress,Mah13,saxena2014progress,saptharishi2019survey}. Here we will formally state a few of the PIT results and properties of circuits for the later reference.

\subsubsection*{Trivial PIT Algorithm}

The simplest PIT algorithm for any circuit in general is due to Polynomial Identity Lemma \cite{ore1922hohere,DL78,Zip79,schwartz1980fast}. When the number of variables is small, say $O(1)$, then this algorithm is very efficient.

\begin{lemma}[Trivial PIT] \label{lem:trivial-pit}
    For a class of $n$-variate, individual degree $ < d$ polynomial $f \in \F[\x]$ there exists a deterministic PIT algorithm which runs in time $O(d^n)$.
\end{lemma}

\subsubsection*{Sparse Polynomial}
Sparse PIT is testing the identity of polynomials with bounded number of monomials. There have been a lot of work on sparse-PIT, interested readers can refer \cite{ben1988deterministic,klivans2001randomness} and references therein. For the proof of poly-time hitting set of Sparse PIT see \cite[Thm.~2.1]{saxena2009progress}.

\begin{theorem}[Sparse-PIT map \cite{klivans2001randomness}] \label{thm:sparse-ks}
Let $p(\x) \in \F[\x]$ with individual degree at most $d$ and sparsity at most $m$. Then, there exists $1 \le r \le (mn \log d)^2$, such that \[ p(y,y^d,\hdots,y^{d^{n-1}}) \ne 0,\bmod\,y^r-1.\] If $p$ is computable by a size-$s$ $\Sigma \Pi$ circuit, then there is a deterministic algorithm to test its identity which runs in time $\poly(s,m)$.
\end{theorem}

Indeed if identity of sparse polynomial can be tested efficiently, product of sparse polynomial can be tested efficiently. We formalise this in the following:

\begin{lemma}[\cite{Sap13} Lemma 2.3] \label{lem:prod-sparse-PIT}
    For a class of $n$-variate, degree $d$ polynomial $f \in \F[\x]$ computable by $\PSP{}{}$ of size $s$, there is a deterministic PIT algorithm which runs in time $\poly(s,d)$.
\end{lemma}

A set $\calH \subseteq \F^n$ is called a Hitting Set for a class polynomial $\calC \subseteq \F[\x]$, if for all $g \in \calC$ \[ g \neq 0 \iff \exists \alphavec \in \calH : g(\alphavec) \neq 0.\]
In literature, PIT has a close association with Hitting set as the two notions are provably equivalent (refer Lemma 3.2.9 and 3.2.10 \cite{Forbes14Thesis}). Note that the set $\calH$ works for every polynomial of the class. Instead of a PIT algorithm occasionally we will use such a set.

\begin{lemma}[Hitting Set of $\PSU{}{}$] \label{lem:prod-sparse-hs}
    For a class of $n$-variate, degree $d$ polynomial $f \in \F[\x]$ computable by $\PSP{}{}$ of size $s$, there is an explicit Hitting Set of size $\poly(s,d)$.
\end{lemma}

\subsubsection*{Algebraic Branching Program (ABP)}
An ABP is a layered directed acyclic graph with $q+1$ many layers of vertices $V_0,\hdots,V_q$ with a source $a$ and a sink $b$ such that all the edges in the graph only go from $a$ to $V_0$, $V_{i-1}$ to $V_i$ for any $i\in [q]$, and $V_q$ to $b$. The edges have {\em uni}variate polynomials as their weights. 
The ABP is said to compute the polynomial \[ f(\x)\;=\;\sum_{p \in \mathsf{paths}(a,b)}\, \prod_{e \in p}\,W(e)\;,\]
where $W(e)$ is the weight of the edge $e$. The ABP has width-$w$ if $|V_i| \le w$, $\forall i \in \{0,\ldots,q\}$. In an equivalent definition, polynomials computed by ABP are of the form $A^T (\prod_{i \in [q]}\, D_i) B$, where $A, B \in \F^{w \times 1}[\x]$, and $D_i \in \F^{w \times w}[\x]$, where entries are univariate polynomials. We encourage interested readers to refer \cite{shpilka2010arithmetic, Mah13} for more detailed discussion.

\begin{definition}[Read-once oblivious ABP (ROABP)]
An ABP is called a {\em read-once oblivious ABP (ROABP)} if the edge weights are univariate polynomials in \emph{distinct} variables across layers. Formally, there is a permutation $\pi$ on the set $[q]$ such that the entries in the $i$-th matrix $D_i$ are univariate polynomials over the variable $x_{\pi(i)}$, i.e.,~they come from the polynomial
ring $\F[x_{\pi(i)}]$.
\end{definition}

A polynomial $f(x)$ is said to be computed by width-$w$ ROABPs in \emph{any order}, if for every permutation $\sigma$ of the variables, there exists a width-$w$ ROABP in the variable order $\sigma$ that computes the polynomial $f(\x)$. In whitebox setting, identity testing of any-order ROABP is completely solved. 

\begin{theorem}[Theorem 2.4 \cite{raz2005deterministic}] \label{thm:white-pit-ro}
    For $n$-variate polynomials computed by size-$s$ ROABP, a hitting set of size $O(s^5 + s \cdot n^4)$ can be constructed.    
\end{theorem}

There have been quite a few results on blackbox PIT for ROABPs as well \cite{forbes2013quasipolynomial,forbes2014hitting,gurjarARO}. The current best known algorithm works in quasipolynomial time.  

\begin{theorem}[Theorem 4.9 \cite{gurjarARO}] \label{thm:roabp-blackbox-pit}
For $n$-variate, individual-degree-$d$ polynomials computed by width-$w$ ROABPs in any order, a hitting set of size $(ndw)^{O(\log \log w)}$ can be constructed.
\end{theorem}

\subsubsection*{Depth-4 Circuits}

A polynomial $f(\x) \in \F[\x]$ is computable by $\SWSP{}{}{[\delta]}$ circuits if $ f(\x) = \sum_{i \in [s]} f_i(\x)^{e_i}$ where $\deg f_i \leq \delta$. The first nontrivial PIT algorithm for this model was designed in \cite{forbes2015deterministic}.

\begin{theorem}[Proposition 4.18 \cite{forbes2015deterministic}] \label{thm:swsp-hs}
There is a $\poly(n,d, \delta \log s)$-explicit hitting set of size $(nd)^{O(\delta \log s)}$ for the class of $n$-variate, degree-($\le d$) polynomials $f(\x)$, computed by $\SWSP{}{}{[\delta]}$-circuit of size $s$.
\end{theorem}

Similarly, $\SWSU{}{}{}$ circuits compute polynomials of the form $f(\x) = \sum_{ i \in [s] } f_i^{e_i}$ where $f_i$ is a sum of univariate polynomials. Using duality trick \cite{saxena2008diagonal} and PIT results from \cite{raz2005deterministic,gurjarARO}, one can design efficient PIT algorithm for $\SWSU{}{}{}$ circuits.

\begin{lemma} [PIT for $\SWSU{}{}$-circuits]\label{lem:swsu-hs}
Let  $P \in \SWSU{}{}$ of size $s$. Then, there exists a $\poly(s)$ (respectively~$s^{O(\log \log s)}$) time whitebox (respectively~blackbox) PIT for the same.
\end{lemma}

\begin{proof}[Proof sketch]
We show that any $g(\x)^e = (g_1(x_1)+\hdots +g_n(x_n))^e$, where $\pdeg(g_i) \le s$ can be written as $ \sum_{j} h_{j1}(x_1) \cdots h_{jn}(x_n)$, for some $h_{j\ell} \in \F[x_{\ell}]$ of degree at most $es$. Define, $G := (y+g_1) \cdots (y+g_n)-y^n$. In its $e$-th power, notice that the leading-coefficient is $\cf_{y^{e(n-1)}}(G^e)= g^e$. So, interpolate on $e(n-1)+1$ many points ($y= \beta_i\in\F$) to get 
$$ \cf_{y^{e(n-1)}}(G^e) \;=\; \sum_{i=1}^{e(n-1)+1}\, \alpha_i\, G^e(\beta_i)\;.$$
Now, expand $G^e(\beta_i) = ((\beta_i +g_1) \cdots (\beta_i+g_n)-\beta_i^n)^e$, by binomial expansion (without expanding the inner $n$-fold product). The top-fanin can be atmost $s \cdot (e+1) \cdot (e(n-1)+1)= O(se^2n)$. The individual degrees of the intermediate univariates can be at most $es$. Thus, it can be computed by an ROABP (of {\em any order}) of size at most $O(s^2e^3n)$. 

Now, if $f=\sum_{j \in [s]} f_j^{e_j}$ is computed by a $\SWSU{}{}{}$ circuit of size $s$, then clearly, $f$ can also be computed by an ROABP (of any order) of size at most $O(s^6)$. So, the whitebox PIT follows from \autoref{thm:white-pit-ro}, while the blackbox PIT follows from Theorem \autoref{thm:roabp-blackbox-pit}.
\end{proof}

Further, $\SWSU{}{}$ can be shown to be closed under multiplication i.e., product of two polynomials, each computable by a $\SWSU{}{}$ circuit, is computable by a single $\SWSU{}{}$ circuit. To prove that we will need an efficient way to write a product of a few powers as a sum of powers, using simple interpolation. For an algebraic proof, see  \cite[Proposition~4.3]{CARLINI20125}. 

\begin{lemma}[Waring Identity for a monomial] \label{lem:waring-rank}
Let $M= x_1^{b_1} \cdots x_k^{b_k}$, where $1 \le b_1 \le \hdots \le b_k$, and {\em roots of unity} $\mathcal{Z}(i):=\{z \in \C : z^{b_i+1}=1\}$. Then, $$M= \sum_{\varepsilon(i) \in \mathcal{Z}(i): i=2,\cdots,k}\, \gamma_{\varepsilon(2),\hdots,\varepsilon(k)} \cdot \left(x_1 + \varepsilon(2)x_2 + \hdots + \varepsilon(k)x_k\right)^d\;,$$ where  $d:= \deg(M)=b_1+\hdots+b_k$, and  $\gamma_{\varepsilon(2),\hdots,\varepsilon(k)}$ are scalars ($\rk(M):=\prod_{i=2}^k\,(b_i+1)$ many).       
\end{lemma}

{\noindent \em Remark.} We actually need not work with $\F=\C$. We can go to a small extension (at most $d^k$), for a monomial of degree $d$, to make sure that $\varepsilon(i)$ exists. 

Using the above lemma we prove the closure result.

\begin{lemma}\label{lem:prod-swsw}
Let $f_i(\x,y) \in \F[y][\x]$, of syntactic degree $\le d_i$, be computed by a $\SWSU{}{}$ circuit of size $s_i$, for $i \in [k]$ (wrt $\x$). Then, $f_1 \cdots f_k$ has $\SWSU{}{}$ circuit of size $O((d_2+1) \cdots (d_k+1)\cdot\,s_1\cdots s_k)$.
\end{lemma}

\begin{proof}
Let $f_i=\sum_{j} f_{ij}^{e_{ij}}$; by assumption $e_{ij} \le d_i$ (by assumption). Then using \autoref{lem:waring-rank}, $f_{1j_1}^{e_{1j_1}} \cdots f_{kj_k}^{e_{kj_k}}$ has size at most $(d_2+1) \cdots (d_k+1) \cdot \left(\sum_{i \in [k]}\,\size(f_{ij_i})\right)$, for indices $j_1,\hdots, j_k$. Summing up for all $s_1 \cdots s_k$ many products (atmost) gives the upper bound.
\end{proof}
